\section{Converting the Forecast into Actions}
\label{sec:actions}

In a production hedge fund environment, time series forecasts are only the first step in a quantitative pipeline. To generate economic value, these forecasts must be translated into actionable portfolio allocations and trading signals. This process involves several critical steps:

\subsection{Signal Generation and Expected Returns}
The point forecasts ($\hat{y}_{t+h}$) generated by our models represent the expected return for a given asset over horizon $h$. In isolation, a positive forecast implies a long signal, while a negative forecast implies a short signal. However, the magnitude of the forecast must cross a transaction-cost threshold to be actionable. The probabilistic forecasts (quantile regression) provide a confidence interval; signals are only generated when the entire interval sits above or below zero, indicating high conviction.

\subsection{Risk-Adjusted Sizing (Kelly Criterion)}
Once a signal is generated, position sizing is determined using risk management frameworks such as the Kelly Criterion or mean-variance optimization. The optimal allocation to a specific series $i$ depends not only on its expected return $\mu_i$ but also on its forecasted volatility $\sigma_i^2$ (which can be derived from the spread of our quantile forecasts):
\begin{equation}
f_i^* = \frac{\mu_i}{\sigma_i^2}
\end{equation}
This ensures that highly volatile series receive smaller capital allocations, even if their expected return is high, maintaining a balanced risk profile across the panel.

\subsection{Portfolio Optimization}
Because the panel contains 36,923 series representing various sub-categories, forecasts cannot be acted upon in isolation. A portfolio optimizer (such as Markowitz Mean-Variance Optimization) takes the vector of expected returns and a covariance matrix to output optimal portfolio weights. The optimizer enforces constraints such as sector neutrality, maximum leverage, and turnover limits.

\subsection{Execution and Slippage}
Finally, target weights are passed to an algorithmic execution engine. The engine minimizes market impact and slippage by breaking large orders into smaller chunks (e.g., using TWAP or VWAP algorithms). The difference between the paper forecast return and the actual realized return after execution costs is continuously monitored to calibrate the forecasting models.

